% Tipo de documento. En este caso es un art�culo, para folios A4, tama�o de la fuente 11pt y con p�gina separada para el t�tulo
\documentclass[a4paper,11pt,titlepage]{article}

% Carga de paquetes necesarios. OrdenesArticle es un paquete personalizado
\usepackage[spanish]{babel} 
%quitado temporalmente -> \RequirePackage[T1]{fontenc}
\RequirePackage[ansinew]{inputenx} 
\usepackage[spanish,cap,cont,title,fancy]{OrdenesArticle}
\usepackage{array}
\usepackage{graphicx}
\usepackage{hyperref}
\usepackage{pifont}
\usepackage{listings}
\usepackage[usenames,dvipsnames]{color}
\usepackage{colortbl}
\usepackage{makeidx}
\hypersetup{bookmarksopen,bookmarksopenlevel=3,linktocpage,colorlinks,urlcolor=blue,citecolor=blue,
						linkcolor=blue,filecolor=blue,pdfnewwindow,
						pdftitle={Pr�ctica Primer Cuatrimestre Ingenier�a Software II},
						pdfauthor={Juan Andrada Romero, Juan Gallardo Casero, Jose Domingo L�pez L�pez},
						pdfsubject={Ingenier�a del Software II}}


% Macro para definir una lista personalizada 
\newenvironment{milista}%
{\begin{list}{\textbullet}%
{\settowidth{\labelwidth}{\textbullet} \setlength{\leftmargin}{\dimexpr\labelsep+\labelwidth+5pt}
\setlength{\itemsep}{\dimexpr 0.5ex plus 0.25ex minus 0.25ex}
\setlength{\parsep}{\itemsep}
\setlength{\partopsep}{\itemsep}
\addtolength{\topsep}{-7.5pt}
}}%
{\end{list}}

% Macro para insertar una imagen
%       Uso: \imagen{nombreFichero}{Factor escala}{Caption (leyenda)}{Label (identificador para referenciarla)}
% -------------------------------------------------------------------------------------------------------------
\def\imagen#1#2#3#4{
 \begin{figure}[h]
 \begin{center}
   \scalebox{#2}{\includegraphics{#1}}
 \caption {#3}
 \label{#4}
 \end{center}
 \end{figure}
}

% Definicion del "listings" para el lenguaje Java
\lstdefinelanguage{Java}
{
 morecomment = [l]{//}, 
 morecomment = [l]{///},
 morecomment = [s]{/*}{*/},
 morestring = [b]", 
 sensitive = true,
 morekeywords = {package, static, while, switch, break, line, void, String, Object, int, Integer, instanceof, else, if, for, private, return, new, public, class, import, int, boolean, true, false, extends, final, super, protected, abstract, this, do, float, double, null, try, catch, implements}
}

% Configuraci�n del c�digo incluido para el lenguaje Java
\lstset{
  language=Java,
  basicstyle=\footnotesize,
  backgroundcolor=\color{white},
  showspaces=false,
  showstringspaces=false,
  showtabs=false,
  frame=single,
  tabsize=2,
  captionpos=b,
  breaklines=true,
  breakatwhitespace=false,
  escapeinside={\%},
  keywordstyle = \color [rgb]{0,0,1},
  commentstyle = \color [rgb]{0.133,0.545,0.133},
  stringstyle = \color [rgb]{0.627,0.126,0.941}
}


\begin{document}

% En las p�ginas de portada e �ndices, no hay encabezado ni pie de p�gina
\pagestyle{empty} 

% Se incluye la portada
\thispagestyle{empty}
\begin{center}
  {\LARGE UNIVERSIDAD DE CASTILLA-LA MANCHA} \\
  \bigskip
  {\Large ESCUELA SUPERIOR DE INFORM�TICA} \\
  \vspace{28mm}
  \includegraphics[scale=0.45, keepaspectratio]{./imagenes/esi_bw.png} \\
  \vspace{30mm}
  {\LARGE \textbf{ Ingenier�a del Software II}} \\
  \vspace{10mm}
  {\large \textsf{\textsc{SSCA - Sistema de Salud de una Comunidad}}} \\
  {\large \textsf{\textsc{Aut�noma}}} \\
  \vspace{20mm}
  {\large Juan Andrada Romero} \\
  {\large Juan Gallardo Casero} \\
  {\large Jose Domingo L�pez L�pez} \\
  \vspace{9mm}
  {\large \today}
\end{center}
\clearpage

% Texto del reverso de la portada
%%\mbox{}
%%\vspace{18cm}
%%\begin{small}
% Se ajusta la separaci�n entre p�rrafos
%%\parskip=10pt 

%%\copyright~ 2008/2009 Juan Andrada Romero. Universidad de Castilla La Mancha, Escuela Superior de Inform�tica de Ciudad Real.

%%Se permite la modificaci�n, copia y distribuci�n de este documento, seg�n la licencia de documentaci�n GNU (\url{http://www.gnu.org}).

%%Este documento fue compuesto con \LaTeX{}. Im�genes generadas con Power Point y Gimp.
%%\end{small}

%%\newpage

% En las p�ginas de �ndices y prefacio, se utiliza numeraci�n romana
\pagenumbering{Roman}

% Se crea el �ndice
\tableofcontents
% Se pasa p�gina y se a�ade el �ndice de figuras al �ndice principal
%\clearpage\phantomsection
%\addcontentsline{toc}{section}{\listfigurename}
% Se crea el �ndice de figuras
%\listoffigures
% Si se quiere crear un �ndice de tablas se pondr�a: \listoftables

\newpage

% Se ajusta la separaci�n entre p�rrafos
\parskip=10pt

% Se a�ade el prefacio al �ndice
\clearpage\phantomsection
\addcontentsline{toc}{section}{Prefacio}

% Comienza el contenido del documento. Se utilizan n�meros ar�bigos y el encabezado y pie de p�gina personalizado
\pagenumbering{arabic}
\pagestyle{fancy}

% Fase de Inicio
\section{Fase de Inicio}

\subsection{Especificaci�n inicial de requisitos} \label{requisitosIniciales}

Se desea desarrollar un sistema para la gesti�n de citas del sistema de salud
de una comunidad aut�noma. En la primera fase del proyecto (correspondiente al
primer cuatrimestre del curso) se debe construir la aplicaci�n que se
utilizar� en las unidades de citaci�n de los hospitales y centros de salud. En
la segunda fase (segundo cuatrimestre) se construir� una aplicaci�n web para
que pueda ser usada por parte de los beneficiarios del sistema de salud. Los
requisitos de esta segunda fase se dar�n oportunamente.

El sistema de salud de la comunidad aut�noma (SSCA) tiene un conjunto muy
amplio de beneficiarios, de cada uno de los cuales se almacenan sus datos
personales habituales (n�mero de seguridad social, NIF o NIE, nombre,
apellidos, domicilio, tel�fono fijo y m�vil, correo electr�nico, etc.).
Cada beneficiario tiene asignado un m�dico de cabecera, que pasa consulta en
un consultorio, centro de salud u hospital. De los m�dicos tambi�n se
almacenan sus datos personales.

Cuando un beneficiario desea pedir cita, se acerca a una unidad de citaci�n o
llama por tel�fono, de manera que le atiende un empleado que es quien que
maneja el sistema que se debe desarrollar en la primera fase. El beneficiario
se identifica con su n�mero de seguridad social, NIF o NIE, y entonces el
empleado ve cu�l es el m�dico de cabecera que tiene asignado. El empleado
propone al beneficiario alguno de los d�as y horas en los que el m�dico le
puede atender. Cuando el beneficiario acepta una cita, el empleado la marca y
�sta le queda reservada.

Cada m�dico puede tener un horario de consultas fijo (por ejemplo, todos los
d�as de 9 a 14) o variable (lunes y martes de 9 a 14, mi�rcoles de 12 a 14 y
de 17 a 20, ...), si bien hay d�as en que no pasan consulta por alg�n motivo
(vacaciones, quir�fano, asuntos propios, baja, etc.). Cuando un m�dico no va
pasar consulta en una fecha, el SSCA designa normalmente un sustituto. Si el
beneficiario del SSCA pide cita y el m�dico va a estar ausente, el empleado le
debe informar del nombre del m�dico que le atender� antes de darle la cita.

El sistema debe permitir la modificaci�n y anulaci�n de citas, as� como dar
citas a beneficiarios que no est�n registrados previamente. El m�dico de
cabecera de los menores de 14 a�os es un pediatra.

En ocasiones, el m�dico de cabecera (o el pediatra) puede enviar al
beneficiario a una consulta con el especialista. En este caso, el m�dico
utiliza el mismo sistema que usted tiene que desarrollar para emitir un
``volante''. El volante permite al beneficiario pedir cita para el
especialista. El m�dico de cabecera solicitar� el volante al sistema, que le
devolver� un identificador �nico de volante. Con este n�mero, el beneficiario
puede ir a la unidad de citaci�n o llamar por tel�fono y pedir cita para el
especialista. El volante relaciona, al menos, al m�dico de cabecera (o pediatra), al m�dico especialista y al beneficiario. Luego, el m�dico
especialista puede expedir m�s volantes para que el beneficiario vuelva a
acudir a su consulta o a la consulta de otro especialista.

Debido a la informaci�n confidencial que se maneja, es necesario que quede un
registro de todas las operaciones realizadas por cada usuario del sistema.

Los empleados son los usuarios del sistema y, por tanto, utilizan los equipos
inform�ticos de las unidades de citaci�n (aplicaciones cliente). Si el usuario
es administrador, adem�s de las operaciones de los usuarios con rol de
``citador'', el usuario puede:
\begin{milista}
\item Crear, modificar y eliminar otros usuarios.
\item Asignar o desasignar roles.
\item Modificar los calendarios y horarios de disponibilidad de los
m�dicos, estableciendo los m�dicos sustitutos cuando sea necesario.
\end{milista}

En modo ``citador'', el cliente mostrar� un calendario con las citas de cada
m�dico en cada d�a. Los datos de usuarios, m�dicos, citas, etc. residen tanto
en la BD principal como en la de respaldo, que son copias id�nticas una de
otra. Cuando se confirma una cita, el servidor A recibe la petici�n de alguno
de los clientes:
\begin{milista}
\item A env�a la petici�n a B
\item B la almacena en su BD de respaldo
\item B env�a a A un mensaje de confirmaci�n
\item A almacena la cita
\item A env�a un mensaje al cliente indic�ndole que la cita se ha dado
\item A actualiza el estado de los clientes: es decir, pasa a mostrar la
hora como ocupada (o como libre, si se trata de una cancelaci�n de cita)
\end{milista}

Si B falla, la transacci�n de que se trate (alta, modificaci�n o lo que sea)
no puede realizarse.

\subsection{Especificaci�n detallada de requisitos}

Para facilitar el an�lisis y el posterior dise�o e implementaci�n del sistema,
desde la fase de inicio del desarrollo su funcionalidad se dividi� en cinco
partes b�sicas, correspondientes a la gesti�n de beneficiarios, usuarios,
citas, volantes y sustituciones. En los pr�ximos apartados se explican los
requisitos detallados de cada parte del sistema.

\subsubsection{Gesti�n de beneficiarios}

El sistema permitir� la creaci�n, consulta, modificaci�n y tambi�n la
eliminaci�n de beneficiarios, a pesar de que en la especificaci�n inicial no
se hac�a referencia a ella. De cada beneficiario, el servidor almacenar� los
siguientes datos:
\begin{milista}
\item \textbf{Datos personales}: NIF, NSS (n�mero de seguridad social), nombre,
apellidos y fecha de nacimiento.
\item \textbf{Datos de contacto}: Direcci�n completa (domicilio, n�mero, piso,
puerta, localidad, provincia y c�digo postal), correo electr�nico, tel�fono
fijo y tel�fono m�vil.
\item \textbf{Datos del sistema}: Centro de salud y m�dico de cabecera o
pediatra asignado.
\end{milista}

Al registrarse en el sistema, el citador pedir� al beneficiario que elija el
centro de salud en el que quiere pasar las citas con su m�dico de cabecera o
pediatra. De esta manera, el beneficiario podr� elegir el centro de salud m�s
cercano a su ubicaci�n, o cualquiera que sea de su preferencia, para ir a ver
al m�dico de cabecera.

Una vez que el beneficiario se haya registrado, el sistema le asignar� un
m�dico de cabecera o pediatra (en funci�n de su edad) del centro seleccionado,
y el citador deber� informarle de cu�l ha sido el m�dico elegido por el
sistema.

El beneficiario podr� pedir al empleado del sistema desarrollado la consulta,
modificaci�n y eliminaci�n de sus datos registrados en el sistema, teni�ndose
que identificar para ello con su NIF o su NSS. En general, todas las
operaciones relacionadas con la gesti�n de beneficiarios las pueden realizar
los usuarios del sistema con rol de ``citador'', no siendo necesaria la
intervenci�n de los administradores. Sin embargo, los m�dicos tambi�n deber�an
poder consultar informaci�n de los beneficiarios, por si necesitaran saber su
fecha de nacimiento o alg�n dato para ponerse en contacto con ellos.

El m�dico de cabecera o pediatra asignado a cada beneficiario no se puede
modificar, salvo si el beneficiario solicita que se le cambie el centro de
salud en el que quiere pasar las citas. En este caso excepcional, que se puede
dar si el beneficiario cambia de direcci�n, el sistema le asignar� un m�dico
diferente del nuevo centro de salud seleccionado. El NIF y el NSS de un
beneficiario es fijo y no se puede modificar en ning�n caso.

Los m�dicos asignados a los beneficiarios de menos de 14 a�os son pediatras,
mientras que los asignados a beneficiarios con 14 o m�s a�os son m�dicos de
cabecera. El sistema, de forma autom�tica, tendr� que cambiar el m�dico
asociado a un beneficiario cuando cumpla los 14 a�os.

\subsubsection{Gesti�n de usuarios}

El sistema permitir� la creaci�n, consulta, modificaci�n y eliminaci�n de 
usuarios. Se considerar� que existen tres roles o tipos de usuarios:
``citadores'', administradores y m�dicos, pudiendo haber m�s de un
administrador en el sistema. De cada uno de estos usuarios, el servidor
almacenar� los siguientes datos:
\begin{milista}
\item \textbf{Datos personales}: NIF, nombre y apellidos.
\item \textbf{Datos de contacto}: Correo electr�nico, tel�fono fijo y tel�fono
m�vil.
\item \textbf{Datos del sistema}: Nombre de usuario, contrase�a y centro de
salud asignado.
\end{milista}

Adem�s, cada m�dico tiene un horario o calendario de trabajo en el que se
define a qu� horas puede pasar consulta en cada d�a de la semana (de lunes a
viernes). Por simplicidad, y debido a que no supone una limitaci�n importante,
cada horario se define con intervalos de una hora, de forma que un m�dico puede
pasar consulta de 10:00 a 14:00, o de 10:00 a 11:00 y de 12:00 a 14:00, pero
no de 10:00 a 13:30, o de 10:15 a 13:15.

La jornada laboral v�lida para los m�dicos se considera que empieza a las 9:00
de la ma�ana y finaliza a las 21:00 de la tarde. Para permitir el registro de
m�dicos de diferentes tipos y jornadas, no se establecer� una limitaci�n
inferior ni superior en el n�mero de horas de trabajo diarias ni semanales de
cada m�dico.

Los m�dicos podr�n ser de tres clases: m�dicos de cabecera, pediatras o
especialistas. De cada especialista se guardar�n, adem�s de sus datos
habituales, su especialidad m�dica, que se utilizar� para facilitar la emisi�n
de volantes a los beneficiarios.

Cuando se registre un nuevo usuario en el sistema, independientemente de su
tipo, se le asignar� de forma autom�tica el centro de salud en el que deber�
trabajar. Este centro, del que el nuevo usuario deber� ser informado, no se
podr� modificar en ning�n caso.

Por simplicidad, cada usuario s�lo podr� tener asociado un rol (y, en el caso
de los m�dicos, una clase de m�dico), que no se podr� modificar una vez que
el usuario se ha registrado, puesto que es improbable que un m�dico de cabecera
pase a ser administrador del sistema o viceversa. El �nico cambio de rol que
se deber� poder hacer es pasar el rol de un ``citador'' a administrador y
viceversa, lo cual es m�s probable que ocurra en la realidad. Por otro lado,
el NIF de un usuario es fijo y tampoco se puede cambiar.

Los usuarios podr�n pedir al administrador del sistema desarrollado la
consulta, modificaci�n y eliminaci�n de sus datos registrados en el sistema,
teni�ndose que identificar para ello con su NIF. De los m�dicos tambi�n se
podr�n cambiar, si es necesario, su horario de trabajo, para a�adir o eliminar
horas. En general, debido a su relevancia, todas las operaciones relacionadas
con la gesti�n de usuarios s�lo las pueden realizar los usuarios del sistema
con rol de administrador.

\subsubsection{Gesti�n de citas}

El sistema permitir� el registro, consulta y anulaci�n de citas a
beneficiarios. Puesto que la modificaci�n de una cita no es m�s que una
anulaci�n de la cita antigua y el registro de la nueva, y dado que el proceso
de tramitaci�n de una cita es mucho menos costoso que el de un beneficiario o
usuario, se acord� que no era necesario que el sistema permitiera la
modificaci�n directa de las citas.

Para facilitar la gesti�n y tramitaci�n de las citas, adem�s de considerar que
el horario de consultas de un m�dico se encuentra dividido en intervalos de
una hora (como ya se coment� en el anterior apartado de gesti�n de usuarios),
las citas tendr�n una duraci�n fija de 10 minutos cada una, y s�lo se podr�n
pedir en intervalos de 10 minutos. Esto significa que se podr�n pedir cita a
las 14:10 o las 19:40, pero no a las 16:22 o 20:03.

Cada m�dico podr� tener, como m�ximo, una cita asignada para un mismo d�a y
hora, mientras que, en principio, cada beneficiario podr� tener varias citas
el mismo d�a a la misma hora. No obstante, los beneficiarios s�lo podr�n pedir
cita para el m�dico de cabecera o pediatra que le asign� el sistema, o para un
especialista si disponen de un volante, por lo que, en condiciones normales,
cada beneficiario s�lo podr� pedir una �nica cita en una misma hora.

Cuando el empleado del sistema vaya a registrar una cita para un beneficiario,
se le debe mostrar a priori qui�n es el m�dico que va a hacer realmente la
consulta, teniendo en cuenta las sustituciones que le pudieran estar haciendo
al m�dico esperado por el beneficiario. Por otro lado, si el beneficiario que
quiere pedir cita no se encuentra registrado, el sistema deber� ofrecer la
posibilidad de registrar el beneficiario siguiendo el proceso que se explic�
en el apartado de gesti�n de beneficiarios.

La unidad de citaci�n desarrollada deber� permitir la consulta de las citas
registradas tanto para un cierto beneficiario como para un cierto m�dico del
sistema. Todas las operaciones relacionadas con la gesti�n de citas las podr�n
llevar a cabo de la misma forma los usuarios con rol de ``citador'' y
administrador.

\subsubsection{Gesti�n de volantes}

El sistema permitir� �nicamente la solicitud de nuevos volantes, no siendo
posible la modificaci�n, anulaci�n o consulta de volantes existentes (salvo
cuando un beneficiario va a pedir una cita utilizando un volante).

Cuando un m�dico desee emitir un volante, deber� identificar al beneficiario
que podr� utilizarlo y, a continuaci�n, el sistema le mostrar� una lista de
los m�dicos especialistas registrados en el sistema para que seleccione el
receptor del nuevo volante. Para facilitar la selecci�n del especialista, los
m�dicos se tendr�n que agrupar seg�n su especialidad m�dica indicada al
registrarse.

Antes de emitir un volante, el m�dico deber� informar al beneficiario del
centro de salud en el que pasa consulta el especialista, debido a que, al
contrario de lo que ocurre con el m�dico de cabecera o pediatra, los
especialistas a los que se env�an los beneficiarios no tienen por qu� trabajar
en el mismo centro de salud que indicaron al registrarse.

Cada volante s�lo se podr� utilizar una �nica vez para pedir cita con el
especialista, aunque si, por cualquier motivo, se anula la cita con la que un
beneficiario utiliz� un volante, podr� usarlo de nuevo para pedir cita con el
especialista en otro momento. No obstante, en cualquier caso todos los
volantes tienen una fecha de caducidad de 6 meses desde el momento de su
emisi�n, y no se podr�n utilizar a partir de esa fecha (aunque se puedan usar
para pedir una cita en un fecha posterior a la caducidad del volante).

La emisi�n de volantes es responsabilidad exclusiva de los m�dicos del sistema,
por lo que ni los citadores ni los administradores del sistema tendr�n permiso
para solicitar nuevos volantes para los beneficiarios.

\subsubsection{Gesti�n de sustituciones}

El sistema permitir� s�lo la asignaci�n de nuevas sustituciones, ya que entre
los requisitos iniciales del sistema no se encuentra la modificaci�n, anulaci�n
o consulta de las sustituciones existentes.

De la misma forma que las horas de consulta de los m�dicos se dividen en
intervalos de una hora, las sustituciones s�lo se podr�n hacer en horas
completas, pudi�ndose solicitar una sustituci�n en varias horas consecutivas o
separadas para un mismo m�dico en un mismo d�a.

Cuando un m�dico est� ausente, su horario de consultas seguir� siendo el mismo,
pero las citas las pasar� el m�dico sustituto. Por tanto, las horas que se
deber�n sustituir ser�n �nicamente aquellas en las que el m�dico sustituido
tendr�a consulta. Un m�dico no podr� sustituir a dos m�dicos diferentes en
las mismas horas, pero s� podr� sustituir a varios m�dicos el mismo d�a si no
se solapan sus horarios de consulta. Por ejemplo, un m�dico que no trabaje los
viernes podr�a sustituir en un viernes a la vez a un m�dico que trabaja de
9:00 a 14:00 y a otro que trabaja de 16:00 a 20:00.

Los administradores ser�n los �nicos usuarios que puedan asignar sustitutos a
los m�dicos. Por tanto, cuando un m�dico no vaya a poder pasar consulta en una
fecha, deber� avisar al administrador, el cual identificar� al m�dico en el
sistema y, a continuaci�n, elegir� el d�a y las horas en las que el m�dico va
a estar ausente. Si el m�dico ya estuviera siendo sustituido o estuviera
haciendo una sustituci�n en las horas elegidas, no se podr�a asignar un nuevo
sustituto.

Con el fin de hacer m�s sencilla para el administrador la tarea de asignar
sustitutos a los m�dicos, el sistema deber� generar autom�ticamente una lista
con todos los m�dicos del sistema que puedan hacer una sustituci�n pedida. Un
m�dico podr� ser elegible para realizar una sustituci�n si cumple los
siguientes requisitos:
\begin{milista}
\item Ser exactamente de la misma categor�a que el m�dico sustituido, lo cual
incluye la especialidad m�dica en el caso de que se vaya a sustituir un
especialista.
\item Trabajar en el mismo centro de salud que el m�dico sustituido, para que
el beneficiario no se tenga que desplazar por culpa de la sustituci�n.
\item No tener horario de consultas en alguna de las horas que se tienen que
sustituir.
\item No estar pendiente de realizar una sustituci�n en alguna de las horas que
se van a tener que sustituir, incluso si a�n no tiene ninguna cita asignada en
esas horas.
\end{milista}

\subsubsection{Otras consideraciones}

Para permitir que el sistema se pueda ejecutar sobre un entorno distribuido,
la unidad de citaci�n, el servidor front-end, el servidor de respaldo, la base
de datos principal y la base de datos secundaria deben poder estar en m�quinas
diferentes sin que ello afecte al funcionamiento del sistema.

Todas las operaciones que realicen los usuarios del sistema deben quedar
registradas en la base de datos por si fuera necesaria consultarlas m�s
adelante. Adem�s, el servidor front-end debe mostrar por pantalla los
resultados de las �ltimas operaciones realizadas para que se pueda ver
f�cilmente si se producen errores en el sistema.


\subsection{Lista de requisitos funcionales}

Despu�s de analizar con detalle los requisitos y las consideraciones tenidas
en cuenta para el desarrollo del sistema, en este apartado se listan todos los
requisitos funcionales del sistema.

\subsubsection{Servidor front-end}

\begin{milista}
\item \textbf{RF1}: Identificar y registrar a los clientes.
\item \textbf{RF2}: Obtener los datos de un beneficiario.
\item \textbf{RF3}: Almacenar un nuevo beneficiario.
\item \textbf{RF4}: Modificar los datos de un beneficiario.
\item \textbf{RF5}: Eliminar un beneficiario.
\item \textbf{RF6}: Obtener los datos de un usuario (o m�dico).
\item \textbf{RF7}: Almacenar un nuevo usuario (o m�dico).
\item \textbf{RF8}: Modificar los datos de un usuario (o m�dico).
\item \textbf{RF9}: Eliminar un usuario (o m�dico).
\item \textbf{RF10}: Almacenar sustituciones a m�dicos.
\item \textbf{RF11}: Consultar las citas de un beneficiario.
\item \textbf{RF12}: Consultar las citas de un m�dico.
\item \textbf{RF13}: Almacenar una nueva cita.
\item \textbf{RF14}: Eliminar una cita.
\item \textbf{RF15}: Almacenar un nuevo volante.
\item \textbf{RF16}: Registrar todas las operaciones realizadas por los
usuarios en la BD.
\item \textbf{RF17}: Actualizar el estado de los clientes cuando se produce una
operaci�n en el servidor.
\item \textbf{RF18}: Enviar peticiones al servidor de respaldo para gestionar
la base de datos secundaria.
\end{milista}

\subsubsection{Servidor de respaldo}

\begin{milista}
\item \textbf{RF1}: Recibir peticiones para gestionar la base de datos secundaria.
\end{milista}

\subsubsection{Unidad de citaci�n}

\begin{milista}
\item \textbf{RF1}: Identificar a los usuarios del sistema.
\item \textbf{RF2}: Consultar los datos de un beneficiario.
\item \textbf{RF3}: Registrar un nuevo beneficiario.
\item \textbf{RF4}: Modificar los datos de un beneficiario.
\item \textbf{RF5}: Eliminar un beneficiario.
\item \textbf{RF6}: Consultar los datos de un usuario (o m�dico).
\item \textbf{RF7}: Registrar un nuevo usuario (o m�dico).
\item \textbf{RF8}: Modificar los datos de un usuario (o m�dico).
\item \textbf{RF9}: Eliminar un usuario (o m�dico).
\item \textbf{RF10}: Asignar sustitutos a los m�dicos.
\item \textbf{RF11}: Consultar las citas de un beneficiario.
\item \textbf{RF12}: Consultar las citas de un m�dico.
\item \textbf{RF13}: Registrar una nueva cita.
\item \textbf{RF14}: Anular una cita.
\item \textbf{RF15}: Emitir un volante.
\item \textbf{RF16}: Recibir notificaciones para actualizar el estado.
\end{milista}


\clearpage

% Fase de Transici�n
\section{Fase de Transici�n}

\subsection{Ejemplo de trazabilidad} \label{trazabilidad}

En este apartado se va a mostrar un ejemplo de trazabilidad de un caso de uso,
con el fin de ver c�mo se ha ido desarrollando la funcionalidad del caso de uso
a partir de los diferentes artefactos del Proceso Unificado de Desarrollo, as�
como la coherencia existente entre los diagramas, la documentaci�n y el c�digo
final del sistema.

El caso de uso que se tomar� como referencia para el ejemplo de trazabilidad
ser� \textbf{Tramitar cita}, debido a la importancia que este requisito
funcional tiene en el sistema.

\subsubsection{Identificaci�n}

Seg�n el \diagrama{cliente}{casosUso}{Casos de uso}, el caso de uso
\textbf{Tramitar cita} lo pueden realizar los citadores y, por herencia,
tambi�n los administradores del sistema. Sin embargo, como se explica en la
descripci�n del caso de uso, es necesario haber iniciado sesi�n primero,
momento en el que se determinar� el rol del usuario.

Por este motivo, cuando comienza la aplicaci�n, lo primero que se hace es
mostrar la ventana de login al usuario, que como se puede observar en el
\diagrama{cliente}{an�lisis}{An�lisis - Iniciar sesi�n}, se define en la clase
\emph{JFLogin}.

\texttt{\lstinputlisting[inputencoding=latin1]{codigos/TrazaInicio.java}}

Cuando el usuario introduce sus datos en la ventana, la interfaz gr�fica
verifica que los datos tengan un formato v�lido, es decir, que tanto el nombre
de usuario como la contrase�a tengan �nicamente caracteres alfanum�ricos, y que
adem�s la contrase�a tenga al menos 8 caracteres (por motivos de seguridad).
S�lo en caso de que el usuario y la contrase�a sean v�lidos, se manda la
petici�n de inicio de sesi�n al controlador, para que se reenv�e al servidor.

Con el fin de evitar la anidaci�n de condiciones cuando se necesitan verificar
campos de una ventana (tanto la de identificaci�n como todas las dem�s del
cliente), el manejo de valores incorrectos se hace mediante excepciones, de
forma que si un valor no tiene el formato esperado, se lanza una excepci�n que
altera el flujo normal de ejecuci�n.

Como se puede observar a partir del c�digo, en la ventana de identificaci�n
tambi�n se puede introducir la IP y el puerto en el que se encuentra el
servidor front-end. Sin embargo, estos son detalles de implementaci�n que no
est�n recogidos ni en los requisitos funcionales ni en los diagramas de
an�lisis.

\texttt{\lstinputlisting[inputencoding=latin1]{codigos/TrazaIdentificacion.java}}

El ControladorCliente, cuando recibe la solicitud de inicio de sesi�n, sigue
el flujo indicado en el \diagrama{cliente}{an�lisis}{An�lisis - Iniciar
sesi�n}:
\begin{milista}
\item Establece conexi�n con el servidor front-end mediante la clase
\emph{ProxyServidorFrontend}.
\item Se identifica ante el servidor con el nombre de usuario y la contrase�a
introducidos en la ventana.
\item Guarda la sesi�n devuelta por el servidor para utilizarla m�s adelante.
\item Crea un nuevo \emph{Cliente} (a trav�s de la clase \emph{RemotoCliente},
que es singleton), lo activa para ponerlo a la escucha y lo registra en el
servidor.
\item Finalmente, oculta la ventana de identificaci�n y muestra la ventana
principal, \emph{JFPrincipal}, con las operaciones disponibles al usuario
(de esto se encarga el m�todo \emph{iniciar}).
\end{milista}

La identificaci�n del usuario, que en el diagrama de clases de an�lisis se
muestra con un �nico paso, se realiza en el c�digo con dos debido a que la
interfaz del servidor front-end propuesta e implementada posee dos m�todos
diferentes para identificarse y registrar el cliente en el sistema:
\emph{identificar} y \emph{registrar}. Las clases que intervienen en este
proceso se pueden ver en el \diagrama{cliente}{clases}{Arquitectura}.

\texttt{\lstinputlisting[inputencoding=latin1]{codigos/TrazaConexion.java}}

Suponiendo que el usuario registrado es un citador o un administrador, la
interfaz mostrar�, entre las operaciones disponibles, la correspondiente al
caso de uso \textbf{Tramitar cita}. Por tanto, en este momento se cumple la
primera precondici�n necesaria para poder llevar a cabo este caso de uso.

Conviene decir que el cliente consulta al servidor las operaciones que tiene
disponibles para mostrar los paneles que sean necesarios en la interfaz
gr�fica. Este hecho no se muestra ni en el diagrama de casos de uso ni en los
diagramas de clases de an�lisis porque es un detalle de implementaci�n del
caso de uso \textbf{Iniciar sesi�n}, que pretende hacer m�s extensible tanto
el cliente como el servidor, para que haya que hacer los menores cambios
posibles si un rol pasa a tener m�s o menos operaciones disponibles.

\subsubsection{Introducci�n del NIF/NSS del beneficiario}

La clase de an�lisis \emph{JFPrincipal} mostrada en todos los diagramas de
clases de an�lisis del cliente se corresponde, en la implementaci�n final, con
un gran n�mero de clases JPanel, que dividen la funcionalidad del cliente en
cinco grandes grupos (gesti�n de beneficiarios, usuarios, citas, volantes y
sustituciones) y cada uno de estos grupos a su vez en las funcionalidades
b�sicas (registrar beneficiario, tramitar cita, asignar sustituto, etc.). Para
ver el esquema de clases que se ha utilizado en la interfaz gr�fica se puede
consultar el \diagrama{cliente}{clases}{Ventana Principal}.

En concreto, el panel utilizado en el caso de uso \textbf{Tramitar cita} es
\emph{JPCitaTramitar}, pero puesto que para tramitar una cita primero se debe
identificar el beneficiario que la quiere pedir, este panel utiliza a su vez
las funciones del panel \emph{JPBeneficiarioConsultar}. Esta reutilizaci�n de
c�digo demuestra que se ha utilizado la potencia de la orientaci�n a objetos
en el dise�o de la interfaz gr�fica del cliente, lo cual ha facilitado su
desarrollo.

Por tanto, los primeros pasos del \diagrama{cliente}{an�lisis}{An�lisis -
Tramitar cita}, en los que se obtiene y valida el NIF/NSS del beneficiario que
quiere pedir cita y se comprueba si existe realmente, los realiza el panel
\emph{JPBeneficiarioConsultar}, tal y como se puede ver en el
\diagrama{cliente}{secuencia}{Tramitar cita (tramitaci�n correcta)}.
El siguiente fragmento de c�digo muestra este proceso.

\texttt{\lstinputlisting[inputencoding=latin1]{codigos/TrazaBuscarBeneficiario.java}}

\subsubsection{Visualizaci�n de los d�as libres para pedir cita}

Cuando se cargan los datos del beneficiario que quiere pedir cita, el panel
\emph{JPBeneficiarioConsultar} avisa al panel \emph{JPCitaTramitar} para que
solicite al servidor las fechas en las que se puede pedir la cita para el
beneficiario, que es el siguiente paso en el flujo de eventos del caso de uso
\textbf{Tramitar cita}. En la implementaci�n final, esto se realiza a trav�s
del evento personalizado \emph{beneficiarioBuscado}, como se puede ver al final
del anterior fragmento de c�digo.

Cuando el panel \emph{JPCitaTramitar} recibe la notificaci�n de que se ha
buscado un beneficiario, solicita al servidor, a trav�s de varios mensajes, los
d�as y horas en los que se puede pedir cita, tal y como se muestra en el
\diagrama{cliente}{secuencia}{Tramitar cita (tramitaci�n correcta)}.
Concretamente, con el fin de facilitar el manejo del sistema por parte del
citador o administrador, se piden tres listas al servidor:
\begin{milista}
\item Lista de d�as en los que el m�dico trabaja pero tiene todas las citas
ocupadas (por lo que, a efectos de pedir una nueva cita, se considerar� como
d�a no laboral para el m�dico).
\item Lista de d�as y horas en las que el m�dico ya tiene una cita y no puede
pasar otra consulta.
\item Lista de horas en las que el m�dico pasa consulta cada d�a de la semana.
\end{milista}

Cuando el cliente recibe la informaci�n anterior del servidor, genera un
calendario con los d�as en los que el m�dico del beneficiario puede pasar
consulta y las horas de cada d�a que el m�dico a�n no tiene ocupadas. Este
proceso es totalmente dependiente de los controles usados en la interfaz
gr�fica y no se detallar� en esta documentaci�n.

\texttt{\lstinputlisting[inputencoding=latin1]{codigos/TrazaRecuperarHoras.java}}

\subsubsection{Selecci�n del d�a de la cita}

En este momento del ejemplo de trazabilidad, el usuario dispone de un
calendario con los d�as y las horas en los que el m�dico asociado al
beneficiario puede pasar cita con �l. Sin embargo, no se ha tenido en cuenta
un requisito funcional importante del sistema que s� se menciona en los
diagramas: el citador o administrador debe saber cu�l es el m�dico real que
dar� la cita (teniendo en cuenta las sustituciones) para avisar al
beneficiario.

Esto se debe a que, en el sistema tal y como se ha implementado, se consulta
el m�dico real que dar� la cita cuando el usuario selecciona una cierta hora.
Es decir, el cliente no recibe a priori informaci�n sobre los m�dicos que dar�
cada una de las posibles citas, sino que cuando se selecciona un d�a (mediante
el calendario) y una hora (a trav�s de un ComboBox) en la interfaz gr�fica, se
env�a una petici�n al servidor para saber cu�l es el m�dico real que dar�a la
cita elegida.

\texttt{\lstinputlisting[inputencoding=latin1]{codigos/TrazaMedicoReal.java}}

Una vez que se ha seleccionado en la interfaz el d�a y la hora definitiva de la
cita, contin�a el flujo de ejecuci�n del caso de uso \textbf{Tramitar cita}.
Despu�s de que el panel \emph{JPCitaTramitar} verifique que el d�a y la hora
seleccionados son v�lidos (es decir, que sea una hora laborable y libre
para el m�dico), env�a a la clase \emph{ControladorCliente} la petici�n de
tramitaci�n de cita con los datos seleccionados.

El controlador reenv�a la petici�n al servidor front-end haciendo uso de la
clase \emph{ProxyServidorFrontend}, a�adiendo a los datos recibidos el id de la
sesi�n del citador o administrador, recuperado a partir de la sesi�n que se
almacen� en el controlador cuando se inici� sesi�n. Como se puede observar,
todo este flujo se corresponde perfectamente con el representado en el
\diagrama{cliente}{an�lisis}{An�lisis - Tramitar cita} y el
\diagrama{cliente}{secuencia}{Tramitar cita (tramitaci�n correcta)}.

\texttt{\lstinputlisting[inputencoding=latin1]{codigos/TrazaRegistrarCita.java}}

\subsubsection{Registro de la cita en el servidor}

Cuando el cliente env�a la solicitud de tramitaci�n de cita, empieza el flujo
de eventos del caso de uso \textbf{Tramitar cita} del servidor front-end. Como
se puede ver en el \diagrama{frontend}{an�lisis}{An�lisis - Tramitar cita}, el
cliente se conecta con el sistema a trav�s del objeto remoto
\emph{RemotoServidorFrontend}, que lo �nico que hace es llamar al gestor
principal del servidor, la clase \emph{ServidorFrontend}.

N�tese que, en el diagrama de clases de an�lisis citado, la interacci�n
comienza con la selecci�n del beneficiario que quiere pedir cita, ya que es un
paso imprescindible para solicitar una cita. Sin embargo, como nos estamos
centrando en la trazabilidad del caso de uso \textbf{Tramitar cita}, no hemos
comentado en detalle c�mo se gestiona la recuperaci�n del beneficiario en el
servidor, y hemos pasado directamente al almacenamiento de la cita.

\texttt{\lstinputlisting[inputencoding=latin1]{codigos/TrazaCitaServidor.java}}

A continuaci�n, conforme al diagrama de clases de an�lisis, la clase
\emph{ServidorFrontend} llama a la clase \emph{GestorCitas} con los datos de la
cita que se quiere tramitar. Esta clase, entonces, realiza toda la interacci�n
mostrada en el \diagrama{frontend}{secuencia}{Tramitar cita (tramitaci�n
correcta)} para comprobar la validez de la cita y almacenarla en el
sistema. Muchas de las comprobaciones que se hacen no ser�an necesarias si
el servidor s�lo pudiera ser utilizado por el cliente que se ha desarrollado,
porque el propio cliente impide enviar peticiones inv�lidas al servidor. Sin
embargo, para aumentar la seguridad del servidor, las comprobaciones se repiten
tambi�n en este sistema.

Algunas de las acciones que realiza el gestor de citas son:
\begin{milista}
\item Comprobar que el usuario ha iniciado sesi�n y tiene permisos para
realizar la operaci�n.
\item Comprobar que existe el beneficiario que quiere pedir cita.
\item Comprobar que existe el m�dico asociado al beneficiario.
\item Comprobar que la fecha de la cita es v�lida (posterior a la actual,
dentro del calendario de trabajo del m�dico, etc.).
\item Comprobar que el m�dico no tiene ya una cita a la misma hora y el mismo d�a.
\item Almacenar la cita en la base de datos del sistema.
\end{milista}

Todas las clases de la implementaci�n final del servidor que intervienen en la
gesti�n de citas se pueden consultar en el \diagrama{frontend}{clases}{Gestor
Citas}.

\texttt{\lstinputlisting[inputencoding=latin1]{codigos/TrazaPedirCitaServidor.java}}

Una vez que finaliza la tramitaci�n, como se indica en el escenario del caso
de uso \textbf{Tramitar cita} del servidor, se escribe una entrada en el log
del servidor y se actualizan las ventanas de estado de los servidores (a
trav�s de la clase \emph{GestorConexionesLog}), y se avisa a los clientes
registrados de la tramitaci�n de la nueva cita (mediante el
\emph{GestorSesiones}). La secuencia de llamadas que realiza el
\emph{GestorConexionesLog} para que todas las clases interesadas procesen el
mensaje se puede ver en el \diagrama{frontend}{clases}{Gestor Conexiones Log}.

\texttt{\lstinputlisting[inputencoding=latin1]{codigos/TrazaLogYClientes.java}}

\subsection{Manual de Usuario}

%requisitos
%instalaci�n de programas (java, mysql)
%ejecuci�n de aplicaciones
%howtos

\subsubsection{Requisitos del sistema}

Las tres aplicaciones que forman parte del sistema se han desarrollado bajo el lenguaje de programaci�n Java para hacer posible su portabilidad entre distintos sistemas, siempre y cuando �stos tengan instalada una m�quina virtual de Java. De este modo, las aplicaciones son compatibles tanto en entornos basados en Windows 98/NT/XP/Vista/Seven como en entornos GNU/Linux.

Si observamos el diagrama despliegue (ver \diagrama{sistema}{despliegue}{Despliegue}) se aprecia que el sistema es altamente distribuido ya que se pueden utilizar m�quinas dedicadas a soportar los sistemas gestores de bases de datos, m�quinas separadas para cada uno de los servidores y una o varias m�quinas para cada una de las unidades de citaci�n que se deseen. N�tese que en esta primera versi�n los sistemas gestores de bases de datos deben ser basados en MySQL \cite{mysql}.

\subsubsection{Instalaci�n y configuraci�n de los entornos}
Las m�quinas destinadas a soportar los sistemas gestores de bases de datos deber�n tener instalado como servicio el servidor MySQL y las m�quinas sobre las cuales se ejecutar�n las distintas aplicaciones cliente o servidor deber�n tener instalado el entorno de ejecuci�n de Java. A continuaci�n se dan instrucciones de c�mo instalar y configurar los entornos en sistemas basados en Windows.

\paragraph{Entorno de Java}
Para comprobar si ya dispone de una instalaci�n del entorno de ejecucic�n de Java en el equipo y que est� configurado correctamente, abra una ventana de comandos \textit{(Inicio -\textgreater Ejecutar -\textgreater cmd)} y escriba "`java -version"'. Debe obtener un resultado similar al de la figura \ref{fig:java-version}:

\imagen{./imagenes/usermanual/java-version.png}{0.80}{C�mo averiguar la versi�n de Java}{fig:java-version}

Si aparece un mensaje de error, pruebe lo siguiente: 
\begin{milista}
	\item Compruebe si tiene instalado Java. Para ello compruebe que exista el directorio Java en \textit{C:\textbackslash Archivos de programa} (t�picamente). Si lo est�, config�relo como se indica m�s abajo.
	\item Si no existe, vaya a la siguiente direcci�n de Internet y descargue la �ltima versi�n del  Java SE Development Kit (JDK) \cite{download-jdk}
\end{milista}

Una vez descargado, ejecute el archivo, elija el directorio de instalaci�n y siga los pasos para completar la instalaci�n. Tras instalarlo, debe configurarlo. Para ello vaya a \textit{Inicio} y haga clic derecho en \textit{Mi PC}. En las opciones, elija la pesta�a de \textit{Opciones Avanzadas} y haga clic en \textit{Variables de entorno}. En la nueva ventana que aparece, busque la variable \textit{Path} y a�ada al final \texttt{;C:\textbackslash Archivos de programa\textbackslash Java\textbackslash jdk1.x.x\_xx\textbackslash bin}, donde "`x"' hace referencia a la versi�n que tiene instalada de Java.

Ese es el directorio de instalaci�n por defecto de Java. Si en la instalaci�n eligi� otro diferente, debe introducir el que eligi� al instalarlo.

Tras esto, pruebe de nuevo a ejecutar "`java -versi�n"' en una ventana de comandos.


\paragraph{Servidor MySQL}

Descargue el servidor de MySQL de \cite{download-mysql}. Ejecute el fichero de instalaci�n y siga los pasos que se indican. Recuerde la contrase�a que establezca para el usuario \textbf{root}.

A continuaci�n descargue e instale el toolkit MySQL Workbench \cite{download-workbench} que contiene la aplicaci�n \textit{MySQL Administrator}. Ejecute \textbf{MySQL Administrator} e introduzca los datos de conexi�n al servidor MySQL instalado en el paso anterior (ver \ref{fig:mysql-admin-connection-settings}).

\imagen{./imagenes/usermanual/mysql-administrator1.png}{0.80}{Conexi�n de MySQL Administrator con el servidor MySQL}{fig:mysql-admin-connection-settings}

Si ha introducido correctamente la direcci�n IP del servidor, el nombre del usuario y la contrase�a, se le abrir� una nueva ventana desde la cual podr� configurar cualquier aspecto relacionado con el servidor MySQL. A continuaci�n se realizar�n tres pasos id�nticos sea cual sea el servidor que est� configurando (principal o de respaldo):
\begin{enumerate}
	\item \textbf{Creaci�n de la base de datos}. En la barra de men�s de la parte superior de la aplicaci�n seleccione \textit{Tools -\textgreater{} MySQL Query Browser} y, en la nueva aplicaci�n que se iniciar�, elija \textit{File -\textgreater{} Open Script...}. Se mostrar� un cuadro de di�logo en el que tendr� que abrir el fichero \textbf{install.txt} del servidor que est� configurando. A continuaci�n, pulse el bot�n \textit{Execute} que se encuentra en la barra de herramientas para crear la base de datos.
	\item \textbf{Creaci�n del usuario}. En el men� de la izquierda de la aplicaci�n \textit{MySQL Administrator} dir�jase a la secci�n \textbf{User Administration}. Una vez all�, cree el usuario correspondiente en funci�n del servidor que est� configurando:
	\begin{milista}
		\item Servidor principal. Usuario \textit{iso}; Contrase�a \textit{osi}.
		\item Servidor de respaldo. Usuario \textit{isor}; Contrase�a \textit{rosi}.
	\end{milista}
	\item \textbf{Configuraci�n de los permisos de usuario}. Por �ltimo, debe dotar al usuario que acaba de crear de los permisos necesarios para manipular la base de datos que ha creado en el primer paso. Para ello, mant�ngase en la secci�n \textbf{User Administration} y seleccione el usuario que acaba crear. Observe que aparece una pesta�a con tres columnas. En la primera columna seleccione la base de datos que ha creado en el primer paso. A continuaci�n seleccione los permisos \textit{SELECT}, \textit{INSERT}, \textit{UPDATE} y \textit{DELETE} y mu�valos a la columna central. Por �ltimo, aplique los cambios realizados (ver \ref{fig:mysql-admin-configure-privileges}).
\end{enumerate}

\imagen{./imagenes/usermanual/mysql-administrator2.png}{0.70}{Configuraci�n de los privilegios de usuario}{fig:mysql-admin-configure-privileges}

\subsubsection{Gu�as de uso}

\paragraph{Servidor front-end} \label{manual:frontend}
Al ejecutar el fichero JAR correspondiente al servidor principal, denominado front-end, obtenemos una ventana como la que se muestra en la figura \ref{fig:frontend-ventanaprincipal}. Si la describimos de arriba abajo se aprecian los siguientes elementos:
\begin{milista}
	\item Una barra de men�s (\textit{Archivo}, \textit{Opciones} y \textit{Ayuda}).
	\item Tres botones para poner el servidor a la escucha, pararlo o cerrarlo.
	\item Un �rea de texto donde se va mostrando el log de las acciones que realizan los clientes sobre el servidor.
	\item Una barra de estado donde se muestra informaci�n relativa a las comunicaciones: clientes conectados, datos de conexi�n del servidor, datos de conexi�n del sistema gestor de bases de datos, y datos de conexi�n del servidor de respaldo.
\end{milista}

\imagen{./imagenes/usermanual/frontend-ventanaprincipal.png}{0.70}{Ventana principal del servidor front-end}{fig:frontend-ventanaprincipal}

Una vez descrita la interfaz, configuraremos el servidor. Haga clic en \textit{Opciones} -\textgreater \textit{Configurar...} y obtendr� una ventana similar a la mostrada en la figura \ref{fig:frontend-configuracion}. En dicha ventana deber� introducir los datos de conexi�n del SGBD principal, los datos de conexi�n del servidor de respaldo (en caso de querer activarlo), y el puerto en el que el servidor front-end deber� ponerse a la escucha. Para confirmar los cambios, haga clic en \textit{Aceptar}.

\imagen{./imagenes/usermanual/frontend-configuracion.png}{0.70}{Ventana de configuraci�n servidor front-end}{fig:frontend-configuracion}

Por �ltimo, ponga el servidor a la escucha haciendo clic en \textit{Conectar}. Cuando alg�n cliente se conecte aparecer� reflejado en la barra de estado. Si hay clientes conectados y desea desconectar o cerrar el servidor, se le pedir� confirmaci�n para forzar la desconexi�n de los clientes activos.

\paragraph{Servidor de respaldo}
Dado que esta aplicaci�n es similar al servidor principal, pero con funcionalidades reducidas, dir�jase a la secci�n \ref{manual:frontend} para leer su gu�a de uso.

\paragraph{Unidad de citaci�n}
Al ejecutar el fichero JAR correspondiente a la Unidad de Citaci�n del sistema, obtenemos una ventana de inicio de sesi�n. En ella se escribir� el nombre de usuario y la contrase�a. Adem�s, si se hace clic en el bot�n \textit{Avanzado \textgreater \textgreater} aparecer�n las opciones avanzadas (ver Figura \ref{fig:cliente-login}), donde podr� configurar la direcci�n IP y el puerto del servidor principal.

\imagen{./imagenes/usermanual/cliente-login.png}{0.70}{Ventana de inicio de sesi�n}{fig:cliente-login}

Una vez que ha iniciado la sesi�n correctamente, obtendr� una ventana con diferentes pesta�as. Cada pesta�a contiene una  lista de operaciones. Las pesta�as, que corresponden a grupos de operaciones, y las listas de operaciones pueden variar en funci�n del rol que desempe�a el usuario en el sistema. Para ver una lista de los distintos roles y las operaciones que puede desempe�ar cada uno dir�jase a la secci�n \ref{operaciones-rol}.

A continuaci�n se muestran las distintas operaciones que se pueden realizar con la aplicaci�n:
\begin{milista}
	\item \textbf{Registrar un beneficiario} (ver Figura \ref{fig:cliente-registrarbeneficiario}). Para poder realizar operaciones con beneficiarios (p.e. tramitar una cita o emitir un volante) es necesario que previamente haya sido dado de alta en el sistema. Este panel es el que permite realizar dicha acci�n. Como se indica, los campos marcados con * son obligatorios y el resto opcionales. Para rellenar los datos es necesario ce�irse a determinadas normas:
	\begin{milista}
		\item El NIF (N�mero de Identificacion Fiscal) debe estar compuesto por 9 d�gitos (incluyendo el cero del principio) y una letra.
		\item El NSS (N�mero de Seguridad Social) debe constar de 12 d�gitos.
		\item El nombre, apellidos, domicilio, localidad y provincia deben ser cadenas v�lidas.
		\item El c�digo postal debe constar de 5 d�gitos que expresen la provincia y el municipio.
		\item El correo electr�nico debe seguir un formato v�lido.
		\item Los n�meros de telefono fijo y m�vil deben constar de 9 d�gitos y empezar por 9 y 6 respectivamente.
		\item Seleccionar un centro de saludo. La lista de centros de salud es extra�da directamente de la base de datos del servidor.
	\end{milista}
	
	Si en alg�n campo no se sigue el formato especificado se mostrar� el mensaje de error correspondiente permitiendo al usuario que rectifique dicho dato. S�lo se podr� registrar el beneficiario cuando se hayan rellenado al menos todos los campos obligatorios y de forma correcta. El m�dico asignado al beneficiario se hace de forma aleatoria, pero siempre entre m�dicos que pertenecen al centro seleccionado en este formulario.
	
	Cabe destacar que los beneficiarios no son usuarios del sistema y, por consiguiente, no podr�n iniciar sesi�n en �l.
	\imagen{./imagenes/usermanual/cliente-registrarbeneficiario.png}{0.50}{Ventana para registrar un beneficiario}{fig:cliente-registrarbeneficiario}
	
	\item \textbf{Consultar, editar o eliminar un beneficiario} (ver Figura \ref{fig:cliente-consultarbeneficiario}). Las operaciones de consulta, edici�n y supresi�n de un beneficiario �nicamente pueden llevarse acabo sobre beneficiarios que previamente hayan sido dados de alta en el sistema. Para llevarlas acabo, basta con seleccionar el tipo de identificaci�n sobre el cual vamos a realizar la b�squeda (NIF o NSS) y escribir dicho n�mero en la caja de texto habilitada para tal fin. A continuaci�n se presionar� el bot�n \textit{Buscar}. Si el usuario no est� dado de alta en el sistema se mostrar� un mensaje de error; si por el contrario encuentra dado de alta en el sistema, se obtendr� un mensaje de verificaci�n y se mostrar�n sus datos.
	
	En este momento �nicamente es posible visualizar su informaci�n. Para poder editarla o eliminar al beneficiario es necesario marcar la casilla \textit{Habilitar edici�n}. Una vez marcada, se habilitar�n las cajas de texto susceptibles de edici�n, el bot�n \textit{Guardar cambios} y el bot�n \textit{Eliminar beneficiario}, permitiendo al usuario que realice las acciones pertinentes. A la hora de editar los cambios es necesario seguir el mismo tipo de los datos que cuando se registr� al beneficiario. N�tese que hay campos como el NIF y el NSS que no son editables. Si por consecuencia de un error fuese necesario editarlos, se recomienda eliminar al beneficiario y volverlo a crear de nuevo.
	\imagen{./imagenes/usermanual/cliente-consultarbeneficiario.png}{0.50}{Ventana para consultar, editar y eliminar un beneficiario}{fig:cliente-consultarbeneficiario}
	
	\item \textbf{Registrar un usuario (administrador, citador o m�dico)} (ver Figura \ref{fig:cliente-registrarusuario}). El procedimiento para registrar un usuario es similar al de registrar beneficiarios. La diferencia radica en que se debe especificar un nombre de usuario, una contrase�a (y su confirmaci�n) de al menos 8 caracteres alfanum�ricos y el tipo de usuario. Esta informaci�n adicional es necesaria ya que este tipo de usuarios forman parte del conjunto de usuarios que pueden iniciar sesi�n en el sistema. Los distintos tipos de usuario son administradores, citadores y m�dicos y, dentro de los m�dicos, pueden ser de cabecera, pediatras o especialistas. Adem�s, si se trata de un m�dico especialista hay que proporcionar su especialidad. Por �ltimo, en caso de que el usuario sea de tipo m�dico, ser� necesario configurar su horario laboral para que se puedan llevar a cabo los procesos de tramitaci�n de citas y establecimiento de sustitutos.
	
		\imagen{./imagenes/usermanual/cliente-registrarusuario.png}{0.50}{Ventana para registrar un usuario (administrador, citador o m�dico)}{fig:cliente-registrarusuario}
		
	El horario laboral se puede configurar de un modo uniforme para los cinco d�as de la semana (calendario base) o puede variar de un d�a a otro. Como se puede observar en la figura \ref{fig:cliente-calendariolaboral}, las jornadas de trabajo se configuran por horas, pudiendo ser horas consecutivas para formar una jornada intensiva u horas no consecutivas para formar una jornada partida. Adem�s, el bot�n \textit{Propagar selecci�n} nos permite configurar los cinco d�as de la semana tal y como estamos configurando el d�a actual.
	
		\imagen{./imagenes/usermanual/cliente-calendariolaboral.png}{0.50}{Ventana de configuraci�n del calendario laboral}{fig:cliente-calendariolaboral}
	
	\item \textbf{Consultar, editar o eliminar un usuario (administrador, citador o m�dico)} (ver Figura \ref{fig:cliente-consultarusuario}). Este panel nos permite la consulta, edici�n o supresi�n de un usuario. Para buscarlo ser� necesario proporcionar su NIF. Al igual que en el caso del registro, la �nica diferencia que aporta este panel con respecto al de consulta de beneficiarios es la necesidad de indicar el nombre de usuario, la contrase�a, el tipo de usuario y, en caso de tratarse de un m�dico, el calendario laboral.
	
	Tenga en cuenta que si trata de eliminar un m�dico que ha sido asignado a beneficiarios o tiene citas pendientes, de forma autom�tica se asignar� un nuevo m�dico -del mismo centro de salud- a dichos beneficiarios y las citas ser�n eliminadas.
	
	Adem�s, si elimina un usuario que en ese momento tiene una sesi�n activa, dicho usuario ser� desconectado autom�ticamente del sistema como medida de seguridad.
	
	\imagen{./imagenes/usermanual/cliente-consultarusuario.png}{0.50}{Ventana para consultar, editar y eliminar un usuario (administrador, citador o m�dico)}{fig:cliente-consultarusuario}
	
	\item \textbf{Emitir un volante} (ver Figura \ref{fig:cliente-emitirvolante}). Cualquier m�dico puede emitir un volante a un beneficiario para un especialista determinado. Para realizar dicha acci�n, �nicamente deber� introducir el NIF o NSS del beneficiario, que debe haber sido de alta previamente en el sistema, y seleccionar la especialidad requerida. La lista de especialistas se ir� actualizando conforme se vaya seleccionando una especialidad u otra. Por �ltimo, seleccione un especialista y emita el volante. El sistema de devolver� el n�mero de identificaci�n del volante y la fecha de caducidad del mismo. Facilite esta informaci�n al beneficiario para que posteriormente pueda tramitar la cita indicando el n�mero identificador del volante.
	
	\imagen{./imagenes/usermanual/cliente-emitirvolante.png}{0.50}{Ventana para emitir volantes}{fig:cliente-emitirvolante}
	
	\item \textbf{Tramitar una cita} (ver Figura \ref{fig:cliente-tramitarcita}). Para tramitar una cita deber� introducir el NIF o el NSS del beneficiario, que una vez m�s debe haber sido dado de alta previamente en el sistema. A continuaci�n deber� introducir un d�a del calendario en el que su m�dico asignado pase consulta. Para realizar esta acci�n tiene dos opciones: introducir la fecha a mano siguiendo un formato \textit{dd/mm/yyyy} o pulsando sobre el icono del calendario, donde �nicamente aparecer�n habilitados los d�as laborales del m�dico.
	
	Una vez que haya introducido un d�a v�lido se cargar� la lista de horas en las que el m�dico puede atender al paciente. N�tese que las horas de color rojo corresponden a horas que ya han sido reservadas para otras citas y no podr�n volver a seleccionarse. Si la caja de texto \textit{M�dico} refleja un nombre distinto al del m�dico asignado para el beneficiario, significa que en la fecha y hora seleccionada hay programada una sustituci�n y dicho m�dico es quien realmente pasar� la consulta.
	
	\imagen{./imagenes/usermanual/cliente-tramitarcita.png}{0.50}{Ventana para tramitar una cita}{fig:cliente-tramitarcita}
	
	\item \textbf{Tramitar una cita con volante} (ver Figura \ref{fig:cliente-tramitarcitavolante}). El funcionamiento es similar al panel de tramitar una cita normal. La �nica diferencia es que el beneficiario debe facilitar al usuario el n�mero identificador del volante para poder llevar a cabo la operaci�n.
	
	\imagen{./imagenes/usermanual/cliente-tramitarcitavolante.png}{0.50}{Ventana para tramitar una cita con volante}{fig:cliente-tramitarcitavolante}
	
	\item \textbf{Consultar o anular las citas de un beneficiario} (ver Figura \ref{fig:cliente-consultarcitasbeneficiarios}). Del mismo modo que el resto de operaciones que implican beneficiarios, el usuario deber� introducir el NIF o el NSS del beneficiario, que  debe haber sido dado de alta previamente en el sistema. Todas las citas pendientes que tenga el beneficiario ser�n mostradas en forma de tabla. Si adem�s se desean ver las citas pasadas, es decir, el historial de citas del beneficiario, deber� pulsar el bot�n \textit{Ver hist�rico de citas}. Para eliminar una cita pendiente, selecci�nela y pulse el bot�n \textit{Anular cita}.
	
\imagen{./imagenes/usermanual/cliente-consultarcitasbeneficiarios.png}{0.50}{Ventana para consultar y anular citas de un beneficiario}{fig:cliente-consultarcitasbeneficiarios}

	\item \textbf{Consultar las citas de un m�dico} (ver Figura \ref{fig:cliente-consultarcitasmedicos}). Este panel funciona de forma similar al de \textit{Consultar o anular citas de un beneficiario} solo que en esta ocasi�n sirve para consultar las citas de un m�dico. Para poder ejecutar la operaci�n basta con introducir el NIF del m�dico. Por defecto se muestran las citas pendientes. El bot�n \textit{Ver hist�rico de citas} tiene la misma funcionalidad que en el caso anterior.
		\imagen{./imagenes/usermanual/cliente-consultarcitasmedicos.png}{0.50}{Ventana para consultar las citas de un m�dico}{fig:cliente-consultarcitasmedicos}
	
	\item \textbf{Consultar las citas propias del m�dico activo} (ver Figura \ref{fig:cliente-consultarcitaspropias}). La funcionalidad de consultar las citas de un m�dico no est� disponible para ser llevada a cabo por un m�dico. Sin embargo, un m�dico s� puede consultar sus propias citas gracias a este panel. El bot�n \textit{Ver hist�rico de citas} tiene la misma funcionalidad que en los casos anteriores.	
	\imagen{./imagenes/usermanual/cliente-consultarcitaspropias.png}{0.50}{Ventana para consultar las citas propias del m�dico con la sesi�n activa}{fig:cliente-consultarcitaspropias}
	
	\item \textbf{Establecer sustituciones} (ver Figura \ref{fig:cliente-sustitucion}). Es posible establecer periodos en los que un m�dico no pasar� consulta y se le asignar� un m�dico sustituto. El administrador deber� indicar el d�a en el que el m�dico desea ser sustituido as� como el intervalo de horas de la sustituci�n. En este momento el sistema se encargar� de buscar todos los m�dicos que podr�an llevar a cabo la sustituci�n, que son aquellos que trabajen en el mismo centro y no les coincide su horario laboral con el del m�dico que va a ser sustituido. Por �ltimo, el administrador seleccionar� un m�dico sustituto de la lista de posibles sustitutos y validar� la operaci�n.

\imagen{./imagenes/usermanual/cliente-sustitucion.png}{0.50}{Ventana para establecer sustituciones}{fig:cliente-sustitucion}

\end{milista}

\subsection{Valoraci�n de dificultades}

%patron observador con uispec


A lo largo del desarrollo del sistema se han encontrado distintas dificultades en todas y cada una de las fases. No obstante, las m�s destacables se han producido en las fases de implementaci�n y pruebas debido a ciertos detalles que est�n estrechamente ligados con la tecnolog�a empleada. A continuaci�n se enumeran algunos de ellos:

\begin{milista}
	\item \textbf{Java Remote Method Invocation (RMI)}: en una de las iteraciones en las que ya se dispon�a de una peque�a versi�n funcional de las tres aplicaciones que forman el sistema, se detectaron unos fallos en el funcionamiento normal cuando los distintos servidores se activaban y desactivaban una y otra vez. Despu�s de investigar en Internet, se detect� que es necesario realizar las operaciones \textit{exportObject} y \textit{unexportObject} cuando se quiere activar o desactivar la exportaci�n del objeto remoto.
	
	\item \textbf{Interfaces de red activas}: este problema est� �ntimamente relacionado con el anterior consiste en las interfaces de red que tiene activas el sistema sobre el que se ejecuta alguna de las aplicaciones. Como cada interfaz de red tiene asignada una direcci�n IP, al exportar un objeto, RMI coger� aleatoriamente la direcci�n IP en la que se encuentra dicho objeto.
	
	Por ejemplo, supongamos que el sistema sobre el que se ejecuta la aplicaci�n cliente tiene dos interfaces de red: la primera con una direcci�n IP privada y la segunda con una direcci�n IP p�blica. Al exportar el objeto, RMI coger� aleatoriamente una direcci�n IP, supongamos que selecciona la IP privada. Ahora supongamos que en otra m�quina fuera de la red privada se ejecuta la aplicaci�n servidor a la cual se conecta el cliente para establecer una comunicaci�n cliente-servidor, pero el objeto exportado tiene asignada una IP privada que el servidor no conoce, por lo que ser� imposible establecer la comunicaci�n servidor-cliente, es decir, en el sentido contrario.
	
	Para resolver este problema, es necesario establecer en qu� direcci�n IP estar� a la escucha la aplicaci�n. Esto se hace mediante la sentencia \texttt{System.setProperty("java.rmi.server.hostname", ip\_address)}, pero a�n es necesario determinar qu� direcci�n IP de entre todas las posibles se seleccionar� para exportar el objeto. Para ello se ha realizado una peque�a funci�n que recorre todas las interfaces de red activas en el sistema y las ordena, para su selecci�n, seg�n prioridad:
	\begin{enumerate}
		\item IP p�blica.
		\item IP privada.
		\item Loopback o localhost (127.0.0.1)
	\end{enumerate}

	\item \textbf{Pruebas de interfaz gr�fica de usuario con UISpec4J}: como ya se detalla en la subsecci�n \ref{informe-pruebas}, tambi�n se han realizado pruebas de la interfaz gr�fica de usuario de las aplicaciones. Uno de los problemas que han surgido tiene que ver con las ventanas emergentes despu�s de realizar alg�n tipo de acci�n. Por ejemplo, un mensaje de informaci�n que indica que una operaci�n se ha hecho correctamente. La soluci�n para manejar este tipo de situacoines se ha encontrado en la p�gina web de UISpec4J \cite{uispec4j}.
	
	Por �ltimo, tambi�n se tuvieron serios problemas a la hora de probar el patr�n de observador que dota a los clientes de la capacidad de actualizar la informaci�n que muestran por pantalla cuando se realiza una operaci�n en el servidor desde cualquier otro terminal. La dificultad se encuentra en que el objeto remoto que exportan los clientes es \textit{singleton}, por lo que al tratar de simular la ejecuci�n de dos clientes durante el dise�o y ejecuci�n de las pruebas existe un conflicto de condici�n de carrera en el que el segundo cliente que inicia sesi�n se liga al objeto remoto exportado, no siendo posible que el servidor se comunique con el primero de ellos. Para resolverlo, se evita que el segundo cliente de prueba que inicia sesi�n exporte el objeto remoto y se registre en el servidor, permitiendo que el servidor se comunique correctamente con el primero de ellos.
		
\end{milista}

\clearpage

% Secci�n de ap�ndices
% Se a�ade el titulo de "Ap�ndices" al �ndice
\addcontentsline{toc}{section}{Ap�ndices}

%\input{Apendices.tex}

% A�adimos la bibliografia al �ndice
\clearpage\phantomsection
\addcontentsline{toc}{section}{\bibname}

% Bibliograf�a. En este caso se usa BibTeX
%\bibliographystyle{plain}
%\pagestyle{plain} 
%\bibliography{Bibliografia}

\end{document}
